\documentclass[11pt,a4paper]{article}
\usepackage[UTF8]{ctex}
\usepackage{amsmath,amssymb}
\usepackage{booktabs}
\usepackage{geometry}
\geometry{margin=2.4cm}
\usepackage{xcolor}
\usepackage{tcolorbox}
\tcbuselibrary{skins,breakable}
\usepackage{listings}
\usepackage[hidelinks]{hyperref}
\hypersetup{colorlinks=true,linkcolor=blue!55!black,urlcolor=teal!60!black}

% ---------- code style ----------
\definecolor{codebg}{RGB}{246,246,248}
\definecolor{kw}{RGB}{160,32,140}
\definecolor{str}{RGB}{176,30,28}
\definecolor{cmt}{RGB}{95,110,120}
\definecolor{num}{RGB}{40,40,160}
\lstdefinestyle{py}{
  language=Python,
  basicstyle=\ttfamily\footnotesize,
  keywordstyle=\color{kw}\bfseries,
  stringstyle=\color{str},
  commentstyle=\color{cmt}\itshape,
  numberstyle=\tiny\color{gray}, numbers=none,
  showstringspaces=false, breaklines=true,
  frame=single, rulecolor=\color{gray!45},
  backgroundcolor=\color{codebg}, framesep=4pt,
  xleftmargin=6pt, columns=fullflexible, keepspaces=true,
  morekeywords={nn,th,torch,spaces,self},
}
\lstset{style=py}

% ---------- callout boxes ----------
\newtcolorbox{srcbox}[1][源码位置]{colback=blue!4,colframe=blue!55!black,
  fonttitle=\bfseries,title={#1},breakable,left=2mm,right=2mm,top=1mm,bottom=1mm}
\newtcolorbox{keybox}[1][要点]{colback=red!3,colframe=red!60!black,
  fonttitle=\bfseries,title={#1},breakable,left=2mm,right=2mm,top=1mm,bottom=1mm}
\newtcolorbox{tiebox}[1][对应前文理论]{colback=green!4,colframe=green!45!black,
  fonttitle=\bfseries,title={#1},breakable,left=2mm,right=2mm,top=1mm,bottom=1mm}

\title{\textbf{SB3 代码导读：策略网络如何按环境自动搭建}\\[2pt]
\large Gymnasium + Stable-Baselines3 源码走读}
\author{深度强化学习课程 · 配套讲义}
\date{}

\begin{document}
\maketitle

\begin{abstract}
\noindent 这份讲义带你读 \textbf{Stable-Baselines3 (SB3)} 的策略网络是怎么搭的：从你给的环境（动作空间 + 观测空间）出发，SB3 如何\emph{自动}选出特征提取器、概率分布和网络头，把它们拼成一个可训练的策略。我们把代码结构和前面课上推过的理论（$\nabla_\theta\log\pi$、softmax/高斯头、tanh-squash、on/off-policy、rollout/replay buffer）一一对应起来。代码片段为\textbf{简化示意}（抓逻辑，非逐字源码），每节给出真实文件位置，便于你打开仓库对照阅读。版本基准：SB3 2.x（已迁移到 Gymnasium），PyTorch 后端。
\end{abstract}

\section{先建立心智模型：SB3 的“策略”是什么}

\begin{keybox}[一句话]
在 SB3 里，“policy”不是 RL 术语里那个 $\pi$，而是\textbf{一整个网络包}：特征提取器 + 策略/价值网络（+ 目标网络）+ 优化器，全打包在一个 \texttt{policy} 对象里。真正输出动作的只是其中一部分。
\end{keybox}

不管什么算法，前向数据流都是同一个三段式：

\begin{center}
\begin{tcolorbox}[colback=gray!4,colframe=gray!50,width=0.95\linewidth]
\centering
\texttt{obs} $\;\to\;$ \textbf{预处理}（图像归一化 / 离散 obs 转 one-hot）$\;\to\;$
\textbf{特征提取器}（按\emph{观测}空间选）\\[2pt]
$\to\;$ \textbf{net\_arch 全连接干}（latent\_pi, latent\_vf）$\;\to\;$
\textbf{分布头}（按\emph{动作}空间选）$\;\to\;$ \texttt{action}
\end{tcolorbox}
\end{center}

两个“按空间分派”的点是全部魔法所在：

\begin{itemize}
\item \textbf{观测空间} $\Rightarrow$ 决定\emph{特征提取器}：向量 obs 用 \texttt{FlattenExtractor}，图像用 \texttt{NatureCNN}，字典用 \texttt{CombinedExtractor}。
\item \textbf{动作空间} $\Rightarrow$ 决定\emph{分布 + 网络头}：离散用 softmax/Categorical，连续用高斯，等等。
\end{itemize}

\begin{srcbox}
\texttt{MlpPolicy / CnnPolicy / MultiInputPolicy} 只是“按观测空间预选好特征提取器”的别名，三者最终都走同一套策略类。打印 \texttt{print(model.policy)} 能看到完整结构（见 \S\ref{sec:read}）。
\end{srcbox}

\section{Gymnasium 这一侧：环境、空间、交互循环}
\label{sec:gym}

后面几节讲 SB3 怎么\emph{消费}动作/观测空间。但这些空间和交互数据全部来自 \textbf{Gymnasium} 环境（Farama 接手 OpenAI Gym 后的继任者）。SB3 $\ge 2.0$ 用的是 \texttt{gymnasium}（\texttt{import gymnasium as gym}），\textbf{不是}已弃用的老 \texttt{gym}。先看环境这侧给你什么。

\subsection{两个方法走天下：reset 与 step}

\begin{lstlisting}
import gymnasium as gym

env = gym.make("CartPole-v1")
obs, info = env.reset(seed=0)              # -> (observation, info dict)

for _ in range(1000):
    action = env.action_space.sample()     # here: a random policy
    obs, reward, terminated, truncated, info = env.step(action)   # 5-tuple
    if terminated or truncated:            # episode finished -> reset
        obs, info = env.reset()
env.close()
\end{lstlisting}

\begin{itemize}
\item \texttt{reset()} 返回 \texttt{(obs, info)} 二元组（老 gym 只返回 obs）。
\item \texttt{step(action)} 返回 \textbf{5 元组} \texttt{(obs, reward, terminated, truncated, info)}。
\item 还暴露 \texttt{env.observation\_space}、\texttt{env.action\_space}、\texttt{env.render()}、\texttt{env.close()}。
\end{itemize}

五元组逐个存什么：

\begin{center}
\small
\begin{tabular}{c l l p{7.0cm}}
\toprule
\# & \textbf{名字} & \textbf{类型} & \textbf{存什么} \\
\midrule
1 & \texttt{observation} & space 元素 & 下一个状态 $s'$（执行 action 后的新观测） \\
2 & \texttt{reward} & float & 即时奖励 $r_t$，即 $(s,a)\to s'$ 拿到的回报 \\
3 & \texttt{terminated} & bool & 到达 MDP 真终止态（成功/失败/掉坑） \\
4 & \texttt{truncated} & bool & 因外部原因截断（时间上限、人为中止），MDP 没真结束 \\
5 & \texttt{info} & dict & 辅助信息：日志指标、调试量、隐藏变量；向量环境 autoreset 时真末状态放这（\texttt{final\_obs}） \\
\bottomrule
\end{tabular}
\end{center}

\noindent（完全可观时 \texttt{observation} $=s'$；部分可观时它只是 $s'$ 的一个视图，即 POMDP。）

\begin{keybox}[{对应你熟悉的 transition $(s,a,r,s',\mathrm{done})$}]
$s$ \textbf{不在}返回里——它是上一步的 obs（或 \texttt{reset} 的返回），是本次 \texttt{step} 的\emph{输入}；$a$ 是你传进去的 \texttt{action}；$r=$\texttt{reward}；$s'=$\texttt{observation}；$\mathrm{done}=$\texttt{terminated or truncated}。所以存进 buffer 的一条样本通常是 \texttt{(s, a, reward, observation, terminated)}——\textbf{自举掩码用 \texttt{terminated} 而非 done}，\texttt{truncated} 只用来决定要不要 \texttt{reset}（理由见下框）。
\end{keybox}

\begin{tiebox}[terminated vs truncated：为什么 Gym 要拆开（直接关系自举）]
v0.26 起 Gymnasium 把老的 \texttt{done} 拆成两个布尔：
\begin{itemize}
\item \texttt{terminated}：到达 MDP 真正的终止态（成功/失败/掉坑），此后没有未来回报。
\item \texttt{truncated}：因\emph{外部}原因被截断（到达时间上限 \texttt{TimeLimit}、人为中止），MDP 本身\emph{没}结束。
\end{itemize}
为什么关键？自举目标
\[
y=r+\gamma\,(1-\texttt{terminated})\,V(s')
\]
只能用 \texttt{terminated} 把 $V(s')$ 清零；若是 \texttt{truncated}（时间到），$s'$ 其实还有未来价值，\textbf{不该}清零，否则 bootstrapping 目标算错、loss 有偏。这正是前文 DQN/TD 自举里 “是否乘 $(1-\text{done})$” 那个 done 的正确含义——它应当是 \texttt{terminated}，\emph{不是} \texttt{terminated or truncated}。
\end{tiebox}

\subsection{空间 spaces：SB3 分派的全部依据}

\texttt{gymnasium.spaces} 定义观测/动作的类型。\textbf{后面“动作侧”“观测侧”两节里的 \texttt{isinstance} 判的就是这些类}：

\begin{lstlisting}
from gymnasium import spaces
import numpy as np

spaces.Box(low=-1.0, high=1.0, shape=(3,), dtype=np.float32)   # continuous vector
spaces.Discrete(4)                        # one of {0,1,2,3}
spaces.MultiDiscrete([3, 5])              # one of {0..2} AND one of {0..4}
spaces.MultiBinary(6)                     # 6 independent bits
spaces.Dict({"img": spaces.Box(0, 255, (3, 64, 64), np.uint8),
             "vec": spaces.Box(-1, 1, (10,))})                 # composite obs

space = spaces.Box(-1, 1, (3,))
a = space.sample()         # draw a valid element
a in space                 # membership / validity check
\end{lstlisting}

速记映射（空间 $\to$ SB3 自动选什么，详见后两节）：\texttt{Box}$\to$高斯/压缩高斯头；图像 \texttt{Box}$\to$ NatureCNN；\texttt{Discrete}$\to$ Categorical(softmax)；\texttt{MultiDiscrete}$\to$ 多 softmax；\texttt{MultiBinary}$\to$ 伯努利；\texttt{Dict}$\to$ CombinedExtractor。

\subsection{写一个自定义环境}

继承 \texttt{gym.Env}，在 \texttt{\_\_init\_\_} 里定义两个 space，实现 \texttt{reset} 与 \texttt{step}：

\begin{lstlisting}
import gymnasium as gym
from gymnasium import spaces
import numpy as np

class GoLeftEnv(gym.Env):
    """Toy: a 1-D line of length N; goal is to reach the left end (state 0)."""
    metadata = {"render_modes": []}

    def __init__(self, size=10):
        super().__init__()
        self.size = size
        self.observation_space = spaces.Box(0, size, shape=(1,), dtype=np.float32)
        self.action_space = spaces.Discrete(2)           # 0 = left, 1 = right

    def reset(self, *, seed=None, options=None):
        super().reset(seed=seed)                          # seeds self.np_random
        self.pos = self.size
        return np.array([self.pos], dtype=np.float32), {}

    def step(self, action):
        self.pos = np.clip(self.pos + (-1 if action == 0 else 1), 0, self.size)
        terminated = bool(self.pos == 0)                  # reached the goal
        truncated = False                                 # TimeLimit wrapper may set this
        reward = 1.0 if terminated else -0.01
        return np.array([self.pos], dtype=np.float32), reward, terminated, truncated, {}
\end{lstlisting}

\begin{keybox}[先用 SB3 的检查器验收]
写完务必跑 \texttt{stable\_baselines3.common.env\_checker.check\_env}，它校验 space 类型、reset/step 返回格式、dtype、5 元组是否合规——比训练崩了再 debug 省事得多。
\end{keybox}

\begin{lstlisting}
from stable_baselines3.common.env_checker import check_env
from stable_baselines3 import PPO

env = GoLeftEnv()
check_env(env)                        # raises if the env is non-conformant
model = PPO("MlpPolicy", env).learn(10_000)

# register to create by id (max_episode_steps -> auto TimeLimit, i.e. sets truncated)
gym.register(id="GoLeft-v0", entry_point=GoLeftEnv, max_episode_steps=100)
env = gym.make("GoLeft-v0")
\end{lstlisting}

\subsection{向量化与 SB3 的接管}

RL 要海量样本，单环境太慢，所以并行跑多个副本、一次 step 一批：

\begin{lstlisting}
# Gymnasium native:
envs = gym.make_vec("CartPole-v1", num_envs=8, vectorization_mode="sync")  # or "async"
obs, info = envs.reset(seed=0)                  # obs shape: (8, 4)
obs, rewards, terminated, truncated, infos = envs.step(envs.action_space.sample())
\end{lstlisting}

\begin{keybox}[autoreset：批里某个子环境结束了怎么办]
向量环境会\textbf{自动重置}结束的子环境（默认 next-step 模式）。结束那一步真正的末状态放在 \texttt{infos["final\_obs"]} 里——\textbf{别}拿被重置后的首状态当 $s'$ 去算自举目标，要用 \texttt{final\_obs}（与上面 \texttt{terminated} 的道理一脉相承）。
\end{keybox}

SB3 自带一套 \texttt{VecEnv}（为兼容历史，\texttt{step} 返回 \texttt{(obs, rewards, dones, infos)}，\texttt{done = terminated or truncated}，末状态放 \texttt{info["terminal\_observation"]}）：

\begin{lstlisting}
from stable_baselines3.common.env_util import make_vec_env
from stable_baselines3.common.vec_env import SubprocVecEnv, VecNormalize

venv = make_vec_env("CartPole-v1", n_envs=8)                              # DummyVecEnv by default
venv = make_vec_env("CartPole-v1", n_envs=8, vec_env_cls=SubprocVecEnv)   # true multiprocessing
venv = VecNormalize(venv, norm_obs=True, norm_reward=True)                # running normalization
\end{lstlisting}

\begin{tiebox}
\texttt{DummyVecEnv}（同进程串行）vs \texttt{SubprocVecEnv}（多进程真并行）对应前文“向量化环境”那节：CPU 仿真是瓶颈时用 \texttt{SubprocVecEnv}。给 SB3 传字符串或单个 env 时，它会\textbf{自动}套上 \texttt{Monitor} + \texttt{DummyVecEnv}，所以你拿到的 obs 永远带 batch 维 \texttt{(n\_envs, *obs\_shape)}。
\end{tiebox}

\subsection{Wrappers：不改环境就改输入输出}

\texttt{gymnasium.Wrapper} 及三个子类按需改造：\texttt{ObservationWrapper / RewardWrapper / ActionWrapper} 分别改观测/奖励/动作。常用现成件：\texttt{TimeLimit}（加 \texttt{truncated}）、\texttt{RecordEpisodeStatistics}、\texttt{FlattenObservation}、\texttt{FrameStack}、\texttt{AtariPreprocessing}（灰度/缩放/跳帧）。SB3 侧还有 \texttt{VecNormalize}、\texttt{VecFrameStack}、\texttt{VecTransposeImage}（HWC→CHW）。

\subsection{端到端：一句话把两侧连起来}

\begin{lstlisting}
import gymnasium as gym
from stable_baselines3 import PPO

# Gym side: defines obs/action spaces, reward, terminated/truncated
env = gym.make("LunarLander-v3")             # Box(8) obs, Discrete(4) action
# SB3 side: reads those spaces -> Flatten extractor + Categorical head + value head
model = PPO("MlpPolicy", env, verbose=1)
model.learn(200_000)
\end{lstlisting}

\begin{keybox}[两侧分工]
\textbf{Gym/Gymnasium} 负责\emph{环境}：定义 observation/action space、在 \texttt{step} 里给出 \texttt{reward} 与 \texttt{terminated/truncated}。\textbf{SB3} 负责\emph{智能体}：读这两个 space 自动搭网络（见后两节），跑交互循环、填 buffer、更新参数。你上新问题时，绝大多数工作量在\textbf{把问题表达成一个合规的 \texttt{gym.Env}}（设计 state/action/reward），网络那侧 SB3 基本自动。
\end{keybox}

\section{动作侧：分布怎么按动作空间分派}

这是和前面理论最直接挂钩的一段。整个分派集中在一个工厂函数里：

\begin{srcbox}
\texttt{stable\_baselines3/common/distributions.py} → \texttt{make\_proba\_distribution()}
\end{srcbox}

\begin{lstlisting}
# simplified: action_space -> probability distribution
from gymnasium import spaces

def make_proba_distribution(action_space, use_sde=False, dist_kwargs=None):
    dist_kwargs = dist_kwargs or {}
    if isinstance(action_space, spaces.Box):            # continuous
        cls = StateDependentNoiseDistribution if use_sde else DiagGaussianDistribution
        return cls(get_action_dim(action_space), **dist_kwargs)
    elif isinstance(action_space, spaces.Discrete):     # one categorical choice
        return CategoricalDistribution(int(action_space.n), **dist_kwargs)
    elif isinstance(action_space, spaces.MultiDiscrete): # several independent categoricals
        return MultiCategoricalDistribution(list(action_space.nvec), **dist_kwargs)
    elif isinstance(action_space, spaces.MultiBinary):   # several independent bits
        return BernoulliDistribution(int(action_space.n), **dist_kwargs)
    else:
        raise NotImplementedError(f"unsupported action space {action_space}")
\end{lstlisting}

\begin{center}
\begin{tabular}{l l l}
\toprule
\textbf{动作空间(Gymnasium)} & \textbf{分布类} & \textbf{对应概率模型} \\
\midrule
\texttt{Discrete(n)} & \texttt{CategoricalDistribution} & softmax over $n$ \\
\texttt{MultiDiscrete} & \texttt{MultiCategoricalDistribution} & 多个独立 softmax \\
\texttt{MultiBinary(n)} & \texttt{BernoulliDistribution} & $n$ 个独立伯努利 \\
\texttt{Box}（连续） & \texttt{DiagGaussianDistribution} & 对角高斯 \\
\texttt{Box} + \texttt{use\_sde=True} & \texttt{StateDependentNoiseDistribution} & gSDE（状态相关噪声） \\
\texttt{Box}（SAC/TD3 内部） & \texttt{SquashedDiagGaussian} & tanh 压缩的高斯 \\
\bottomrule
\end{tabular}
\end{center}

\subsection{离散：Categorical 头 = 一个线性层 + softmax}

每个分布类都有个 \texttt{proba\_distribution\_net()}，负责\emph{造网络头}（把 latent 映射成分布参数），另有 \texttt{sample/mode/log\_prob/entropy} 负责采样与梯度：

\begin{lstlisting}
class CategoricalDistribution(Distribution):
    def __init__(self, action_dim):
        self.action_dim = action_dim
    def proba_distribution_net(self, latent_dim):
        # the "softmax head": latent -> one logit per action
        return nn.Linear(latent_dim, self.action_dim)
    def proba_distribution(self, action_logits):
        self.distribution = th.distributions.Categorical(logits=action_logits)
        return self
    def log_prob(self, actions): return self.distribution.log_prob(actions)
    def entropy(self):           return self.distribution.entropy()
    def sample(self):            return self.distribution.sample()
    def mode(self):              return th.argmax(self.distribution.probs, dim=1)
\end{lstlisting}

\begin{tiebox}
注意 \texttt{nn.Linear} 出的是 \emph{logits}，softmax 和 $\log\pi$、熵、采样全交给 PyTorch 的 \texttt{th.distributions.Categorical} —— $\nabla_\theta\log\pi(a|s)$ 由 autograd 自动算，正是前文那个 score。课上推的 $\partial\log\pi(a)/\partial z_b=\mathbb{1}[a{=}b]-\pi(b)$ 就藏在 \texttt{Categorical.log\_prob} 的反向里。
\end{tiebox}

\subsection{连续：高斯头 = 线性均值 + 一个独立的 log\_std 参数}

\begin{lstlisting}
class DiagGaussianDistribution(Distribution):
    def __init__(self, action_dim):
        self.action_dim = action_dim
    def proba_distribution_net(self, latent_dim, log_std_init=0.0):
        mean_actions = nn.Linear(latent_dim, self.action_dim)              # state-DEPENDENT mean
        log_std = nn.Parameter(th.ones(self.action_dim) * log_std_init)    # state-INDEPENDENT std!
        return mean_actions, log_std
    def proba_distribution(self, mean_actions, log_std):
        self.distribution = th.distributions.Normal(mean_actions, log_std.exp())
        return self
    def log_prob(self, actions):
        # independent dims -> sum log-probs over action dimensions
        return sum_independent_dims(self.distribution.log_prob(actions))
    def entropy(self):
        return sum_independent_dims(self.distribution.entropy())
    def sample(self):
        # reparameterized sample: mu + sigma * eps
        return self.distribution.rsample()
\end{lstlisting}

\begin{tiebox}
这正是前文说的“连续策略 = MLP 出 $\mu_\theta(s)$，配一个标准差 $\sigma$”，而且 SB3 默认用\textbf{状态无关}的 $\sigma$：\texttt{log\_std} 是个 \texttt{nn.Parameter}（整段训练共享一个向量），\emph{不}由网络按状态输出。$\nabla_\mu\log\pi=(a-\mu)/\sigma^2$ 同样由 autograd 给出。
\end{tiebox}

\subsection{SAC 的 tanh 压缩高斯：多一个换元修正}

有界连续动作（如力矩 $\in[-1,1]$）不能用裸高斯。SAC 用 \texttt{SquashedDiagGaussianDistribution}：先采高斯 $u$，再 $a=\tanh(u)$，log-prob 要加 tanh 的 Jacobian 修正：

\begin{lstlisting}
class SquashedDiagGaussianDistribution(DiagGaussianDistribution):
    def sample(self):
        self.gaussian_actions = super().sample()          # u ~ N(mu, sigma)
        return th.tanh(self.gaussian_actions)              # a = tanh(u) in (-1, 1)
    def log_prob(self, actions, gaussian_actions=None):
        if gaussian_actions is None:                       # invert tanh if needed
            gaussian_actions = TanhBijector.inverse(actions)
        log_prob = super().log_prob(gaussian_actions)
        # change-of-variables correction: -sum log(1 - tanh(u)^2)
        log_prob -= th.sum(th.log(1 - actions ** 2 + 1e-6), dim=1)
        return log_prob
\end{lstlisting}

\begin{tiebox}
对应前文“连续动作有界 → tanh-squash，log-prob 要加 Jacobian 修正项”。修正项 $-\sum_i\log(1-a_i^2)$ 就是 $\big|\det \partial\tanh/\partial u\big|^{-1}$ 取对数。
\end{tiebox}

\section{观测侧：特征提取器怎么按观测空间选}

\begin{srcbox}
\texttt{stable\_baselines3/common/torch\_layers.py}（\texttt{FlattenExtractor}, \texttt{NatureCNN}, \texttt{CombinedExtractor}, \texttt{MlpExtractor}）；判别图像用 \texttt{common/preprocessing.py} → \texttt{is\_image\_space}
\end{srcbox}

\begin{lstlisting}
class FlattenExtractor(BaseFeaturesExtractor):     # vector / low-dim obs
    def __init__(self, obs_space):
        super().__init__(obs_space, get_flattened_obs_dim(obs_space))
        self.flatten = nn.Flatten()
    def forward(self, obs):
        return self.flatten(obs)                   # the "extractor" is basically a no-op

class NatureCNN(BaseFeaturesExtractor):            # image obs (Mnih et al. 2015)
    def __init__(self, obs_space, features_dim=512):
        super().__init__(obs_space, features_dim)
        n_c = obs_space.shape[0]                    # channels-first (C, H, W)
        self.cnn = nn.Sequential(
            nn.Conv2d(n_c, 32, 8, stride=4), nn.ReLU(),
            nn.Conv2d(32, 64, 4, stride=2), nn.ReLU(),
            nn.Conv2d(64, 64, 3, stride=1), nn.ReLU(), nn.Flatten())
        # one dummy forward to get flattened size, then a Linear to features_dim
        self.linear = nn.Sequential(nn.Linear(self._n_flatten(obs_space), features_dim), nn.ReLU())
    def forward(self, obs):
        return self.linear(self.cnn(obs))
\end{lstlisting}

\begin{itemize}
\item \textbf{1D 向量 obs}（如 CartPole 的 \texttt{Box(4,)}）：提取器就是 \texttt{Flatten}，真正干活的是后面的 \texttt{net\_arch} 全连接。
\item \textbf{图像 obs}（如 Atari 的 \texttt{Box(4,84,84)}）：用 \texttt{NatureCNN}。on-policy（A2C/PPO）默认\textbf{actor/critic 共享} CNN 省算力；off-policy（SAC/TD3）actor 和 critic \textbf{各用一套}提取器（实测效果更好）。
\item \textbf{字典 obs}（\texttt{Dict}，多模态/HER）：用 \texttt{CombinedExtractor}——图像 key 走 CNN，向量 key 直接 flatten，最后拼成一个长向量。
\end{itemize}

\begin{keybox}[预处理悄悄做的事]
喂进提取器前 SB3 会预处理：图像归一化到 $[0,1]$；\textbf{离散观测自动转 one-hot}；向量 obs 仅 flatten。所以你自定义网络拿到的已经是干净的浮点张量。
\end{keybox}

\section{把两侧拼起来：on-policy 的 ActorCriticPolicy}
\label{sec:onpolicy}

A2C / PPO / TRPO 都用 \texttt{ActorCriticPolicy}。它把“特征提取器 + 共享干 + 分布头 + 价值头”组装起来：

\begin{srcbox}
\texttt{stable\_baselines3/common/policies.py} → \texttt{ActorCriticPolicy.\_build / forward / \_get\_action\_dist\_from\_latent / evaluate\_actions}
\end{srcbox}

\begin{lstlisting}
class ActorCriticPolicy(BasePolicy):           # A2C / PPO / TRPO  (simplified)
    def _build(self, lr_schedule):
        self.features_extractor = self.make_features_extractor()       # Flatten / CNN / Combined
        self.mlp_extractor = MlpExtractor(self.features_dim, self.net_arch, self.activation_fn)
        latent_pi = self.mlp_extractor.latent_dim_pi
        # ---- action head: chosen by the distribution type ----
        if isinstance(self.action_dist, DiagGaussianDistribution):
            self.action_net, self.log_std = self.action_dist.proba_distribution_net(
                latent_pi, log_std_init=self.log_std_init)
        elif isinstance(self.action_dist, CategoricalDistribution):
            self.action_net = self.action_dist.proba_distribution_net(latent_pi)
        # (MultiCategorical / Bernoulli / gSDE handled similarly)
        # ---- value head ----
        self.value_net = nn.Linear(self.mlp_extractor.latent_dim_vf, 1)

    def forward(self, obs, deterministic=False):
        features = self.extract_features(obs)
        latent_pi, latent_vf = self.mlp_extractor(features)            # shared / split trunk
        dist = self._get_action_dist_from_latent(latent_pi)
        actions = dist.get_actions(deterministic=deterministic)
        log_prob = dist.log_prob(actions)
        values = self.value_net(latent_vf)
        return actions, values, log_prob          # <-- exactly what goes into the rollout buffer

    def _get_action_dist_from_latent(self, latent_pi):
        out = self.action_net(latent_pi)                               # logits OR mean
        if isinstance(self.action_dist, DiagGaussianDistribution):
            return self.action_dist.proba_distribution(out, self.log_std)
        elif isinstance(self.action_dist, CategoricalDistribution):
            return self.action_dist.proba_distribution(action_logits=out)
\end{lstlisting}

\begin{tiebox}
\texttt{forward()} 返回的 \texttt{(actions, values, log\_prob)} 正是前文 \textbf{rollout buffer} 里存的三样东西。PPO 更新时调 \texttt{evaluate\_actions(obs, actions)} 对 minibatch \emph{重算} \texttt{log\_prob / entropy / values}，再配旧 \texttt{log\_prob} 算重要性比 $\pi_\theta/\pi_{\text{old}}$ —— 这就是前文“近端 off-policy、跑 $K$ 个 epoch”的代码落点。\texttt{net\_arch=dict(pi=[...], vf=[...])} 控制 actor/critic 各自（或共享）的干。
\end{tiebox}

\section{off-policy 三家：没有分布 / 有压缩高斯 / 确定性}

\subsection{DQN：根本没有分布，输出每个动作的 Q}

\begin{srcbox}
\texttt{stable\_baselines3/dqn/policies.py} → \texttt{QNetwork}
\end{srcbox}

\begin{lstlisting}
class QNetwork(BasePolicy):                       # DQN: discrete actions only
    def __init__(self, obs_space, action_space, net_arch, ...):
        self.features_extractor = ...
        # MLP: features -> Q-value for EACH discrete action
        self.q_net = nn.Sequential(*create_mlp(self.features_dim, action_space.n, net_arch))
    def forward(self, obs):
        return self.q_net(self.extract_features(obs))   # shape (batch, n_actions)
    def _predict(self, obs, deterministic=True):
        return self.forward(obs).argmax(dim=1)          # greedy; epsilon handled in DQN.predict()
\end{lstlisting}

\begin{tiebox}
对应前文 DQN 那节：输出 $Q(s,\cdot)$ 的一整排值、\texttt{argmax} 选动作，\textbf{只支持 \texttt{Discrete}}；$\epsilon$-greedy 在外层 \texttt{predict} 里加；target network、replay buffer 在 \texttt{DQN} 算法类里维护，不在策略类里。
\end{tiebox}

\subsection{SAC：状态相关 std 的压缩高斯 + 双 Q}

\begin{srcbox}
\texttt{stable\_baselines3/sac/policies.py} → \texttt{Actor}（演员）、\texttt{ContinuousCritic}（评论家）
\end{srcbox}

\begin{lstlisting}
class Actor(BasePolicy):                           # SAC actor: continuous only
    def __init__(self, ...):
        self.latent_pi = nn.Sequential(*create_mlp(self.features_dim, -1, net_arch))  # trunk
        self.mu      = nn.Linear(last_dim, action_dim)
        self.log_std = nn.Linear(last_dim, action_dim)   # state-DEPENDENT std (cf. on-policy Parameter)
        self.action_dist = SquashedDiagGaussianDistribution(action_dim)
    def forward(self, obs, deterministic=False):
        latent = self.latent_pi(self.extract_features(obs))
        mean = self.mu(latent)
        log_std = self.log_std(latent).clamp(LOG_STD_MIN, LOG_STD_MAX)   # numerical safety
        return self.action_dist.actions_from_params(mean, log_std, deterministic)
\end{lstlisting}

两点和 on-policy 不同：（i）\texttt{log\_std} 是\textbf{网络输出、状态相关}的（PPO 是全局 \texttt{Parameter}）；（ii）评论家 \texttt{ContinuousCritic} 是 \textbf{$n$ 个并列 Q 网络}（默认 2 个，twin-Q 抗高估），且 actor/critic 用\emph{独立}特征提取器。

\subsection{TD3/DDPG：确定性演员 + tanh}

演员直接输出动作 $a=\tanh(\text{MLP}(s))$（无分布、确定性），评论家是 Q 网络。和 SAC 共用 \texttt{ContinuousCritic}。

\section{不同环境 → 自动选了什么：速查表}

\begin{center}
\small
\begin{tabular}{l l l l l}
\toprule
\textbf{环境(例)} & \textbf{观测空间} & \textbf{动作空间} & \textbf{特征提取器/分布} & \textbf{常用算法} \\
\midrule
CartPole & \texttt{Box(4,)} & \texttt{Discrete(2)} & Flatten / Categorical & PPO,A2C,DQN \\
Pendulum & \texttt{Box(3,)} & \texttt{Box(1,)} & Flatten / 高斯(或压缩) & PPO,SAC,TD3 \\
LunarLander-cont. & \texttt{Box(8,)} & \texttt{Box(2,)} & Flatten / 高斯 & PPO,SAC \\
Atari (Breakout) & \texttt{Box(4,84,84)} & \texttt{Discrete} & NatureCNN / Categorical & PPO,DQN \\
MultiDiscrete 任务 & 视情况 & \texttt{MultiDiscrete} & */多 softmax & PPO,A2C \\
机器人+目标(HER) & \texttt{Dict} & \texttt{Box} & CombinedExtractor / 压缩高斯 & SAC,TD3 \\
\bottomrule
\end{tabular}
\end{center}

\begin{keybox}[默认网络大小（1D obs）]
PPO/A2C/DQN：两层 $\times 64$；SAC：两层 $\times 256$；TD3/DDPG：\texttt{[400, 300]}（沿用原论文）。图像一律先过 NatureCNN。
\end{keybox}

\begin{tiebox}[联系你的运筹方向]
路由 / 调度 / DVRP 这类\textbf{组合、可变、结构化}动作空间，SB3 原生分布（flat softmax / 高斯）搞不定。实务路线：（a）\texttt{MaskablePPO}（在 \textbf{sb3-contrib} 里）——对不可行动作做 logit 屏蔽，仍是 Categorical 但带 mask；（b）超出 mask 的（pointer/attention 序贯解码）就得自定义策略类，或转用更灵活的框架。$\nabla\log\pi$ 可微可采样这个内核不变，变的只是“头”的结构。
\end{tiebox}

\section{怎么真正打开仓库读：路线与技巧}
\label{sec:read}

\textbf{推荐阅读顺序}（自底向上，每个文件只看上面点到的类/函数）：
\begin{enumerate}
\item \texttt{common/distributions.py} —— 先看 \texttt{make\_proba\_distribution}，再看 \texttt{Categorical/DiagGaussian} 的 \texttt{proba\_distribution\_net}。这是“动作空间 $\to$ 头”的核心。
\item \texttt{common/torch\_layers.py} —— \texttt{FlattenExtractor / NatureCNN / MlpExtractor / create\_mlp}。这是“观测空间 $\to$ 特征 + 全连接干”。
\item \texttt{common/policies.py} —— \texttt{BasePolicy} 与 \texttt{ActorCriticPolicy}，看 \texttt{\_build} 怎么把上面两块拼起来、\texttt{forward} 怎么产出 \texttt{(action, value, log\_prob)}。
\item 再按需读 \texttt{dqn/policies.py}、\texttt{sac/policies.py}、\texttt{td3/policies.py} 看 off-policy 的差异。
\end{enumerate}

\textbf{最快的“看清结构”技巧}——直接打印策略对象，把上面所有抽象变成具体层：

\begin{lstlisting}
import gymnasium as gym
from stable_baselines3 import PPO

m = PPO("MlpPolicy", "CartPole-v1")
print(m.policy)                 # whole bundle: features_extractor + mlp_extractor + action_net + value_net
print(m.policy.action_dist)     # CategoricalDistribution(...)
print(m.policy.action_net)      # Linear(in_features=64, out_features=2)   <- the softmax head

m2 = PPO("MlpPolicy", "Pendulum-v1")
print(m2.policy.action_dist)    # DiagGaussianDistribution(...)
print(m2.policy.action_net)     # Linear(in_features=64, out_features=1)   <- the mean head
print(m2.policy.log_std)        # Parameter of shape (1,)  <- state-independent std
\end{lstlisting}

\begin{srcbox}[在线源码]
逐字源码可在 readthedocs 的 “\,[source]\,” 链接或 GitHub \texttt{DLR-RM/stable-baselines3} 看。建议对照本讲义的简化版读真实实现，注意它多出来的初始化、类型注解、设备/数据类型处理等工程细节——逻辑骨架与这里一致。
\end{srcbox}

\section{一页纸小结}

\begin{keybox}[把整条链记住]
环境给你两个空间。\textbf{观测空间}选特征提取器（向量 Flatten / 图像 CNN / 字典 Combined）；\textbf{动作空间}经 \texttt{make\_proba\_distribution} 选分布与头（离散 softmax / 连续高斯 / SAC 压缩高斯 / DQN 无分布直接 Q）。中间的 \texttt{net\_arch} 全连接干把特征变成 \texttt{latent\_pi, latent\_vf}。on-policy 的 \texttt{ActorCriticPolicy.forward} 吐出 \texttt{(action, value, log\_prob)} 灌进 rollout buffer；off-policy 各算法在自己的策略类里搭 actor/critic（+target）配 replay buffer。前文所有理论（score、高斯/softmax 头、tanh-squash、on/off-policy、两类 buffer）在这套代码里都能一一对上号。
\end{keybox}

\end{document}
